\section{Funktionen}
\subsection{Default Argumente}
\begin{itemize}
    \item \textbf{Prototipo}: I valori predefiniti si assegnano \textbf{solo} nell'interfaccia.
    \item \textbf{Omissione}: I parametri con default possono essere saltati nella chiamata.
    \item \textbf{Ordine}: I parametri di default devono trovarsi \textbf{\cor{sempre alla fine (a destra)}} della lista parametri.
\end{itemize}

\lstinputlisting[language=c++]{code/02_funktionen/default_argumente.cpp}

Default Argumente können überall verwendet werden bis auf Aufrufparameter von Operatoroverloading. 
Dort sind sie \textbf{nicht} erlaubt. 
Dort machen sie aber ohnehin keinen Sinn. 

\subsection{Funzioni Inline}
\begin{itemize}
    \item \textbf{Sostituto Macro}: Alternativa sicura alle macro C.
    \item \textbf{Zero Overhead}: Codice inserito direttamente (nessun salto di funzione).
    \item \textbf{Type Safety}: Il compilatore garantisce il controllo dei tipi.
    \item \textbf{Target}: Funzioni piccole chiamate frequentemente (es. cicli).
    \item \textbf{No-go}: Inapplicabile a ricorsione o puntatori a funzione.
    \item \textbf{Header-only}: Se usato nei file Header, definizione func. direttamente in file.h
    \item \textbf{Ottimizzazione}: Richiede flag attivi (es. \texttt{-O3}).
\end{itemize}

\lstinputlisting[language=C++]{code/02_funktionen/02_funktionen_snippet_1.cpp}